\section{Eleventh to Thirteenth Centuries: Sacrament, Dualism, Poverty, and Schools}

The word \emph{heresy} now operated within a denser system of canon law, papal government, universities, and secular enforcement. Some groups left their own books; others survive only as names in legislation or portraits by missionaries and inquisitors. The repeated medieval machinery of suppression does not itself prove that every named group held one theology.

\begin{heresydossier}{Berengar of Tours and Eucharistic conversion}{dossier:berengar}
\objectfield Berengar of Tours questioned accounts of Eucharistic conversion in the mid-eleventh century. His surviving works, opponents such as Lanfranc, and several professions document a controversy whose vocabulary changed across repeated proceedings. Evidence grade A/B.

\propositionfield Berengar resisted a substantial conversion that leaves only the appearances of bread and wine and emphasized sacramental signification. Opponents charged him with leaving ordinary bread and wine after consecration. His precise position developed, so one recantation formula cannot simply be treated as his lifelong creed.

\responsefield Roman and regional councils from 1050 onward condemned his teaching. The Roman profession of 1059 used strongly physical language; under Gregory VII, the Roman council of 1079 required a more carefully framed profession that bread and wine are substantially converted into Christ's true body and blood. Lateran IV constitution 1 in 1215 supplied the durable conciliar term: bread and wine are transubstantiated by divine power.

\aftermathfield Berengar submitted, again disputed interpretation, and finally reconciled. The supersession of the 1059 formula by the 1079 settlement cautions against presenting every demanded profession as an equally exact timeless definition. The controversy helped refine sacramental realism while preserving sacramental mode from crude cannibalistic interpretation.
\end{heresydossier}

\begin{heresydossier}{The rupture of 1054 and the limits of a heresy label}{dossier:east-west-1054}
\objectfield Centuries of linguistic, political, liturgical, and jurisdictional estrangement formed the background to the exchange of censures between Cardinal Humbert's legation and Patriarch Michael Cerularius in 1054. The texts survive, but the papal authority of the legates after Leo IX's death and the scope of both sentences require qualification. Evidence grade A/B.

\propositionfield The controversy included Roman primacy, the western addition of the \emph{filioque}, use of leavened or unleavened Eucharistic bread, clerical discipline, and mutual accusations assembled from old heresiological labels. No one proposition or date created a homogeneous ``Orthodox heresy,'' and the legatine bull did not excommunicate every eastern Christian.

\responsefield Humbert placed a bull against Cerularius and specified associates on the altar of Hagia Sophia; a Constantinopolitan synod anathematized the legates. Later reunion councils at Lyons II in 1274 and Florence in 1439 confessed papal primacy and the procession of the Holy Spirit while respecting legitimate eastern rites. The 1054 acts were mutual disciplinary censures within a broader doctrinal and ecclesiological dispute, not an ecumenical dogmatic definition.

\aftermathfield Communion continued locally after 1054 but separation hardened through later events, especially the Latin sack of Constantinople in 1204. Paul VI and Athenagoras removed the 1054 sentences from living ecclesial memory in 1965; they did not declare remaining disagreements solved or erase the history that produced them.
\end{heresydossier}

\begin{heresydossier}{Tanchelm: a dossier preserved by accusers}{dossier:tanchelm}
\objectfield Tanchelm preached in the Low Countries early in the twelfth century and was entangled in a political dispute over the church of Utrecht. A letter of the Utrecht canons to Archbishop Frederick of Cologne, written about 1112--1114, and a later life of Norbert of Xanten are the principal witnesses; no writing by Tanchelm survives. Evidence grade C/D.

\propositionfield The Utrecht letter attributes a Donatist-like dependence of sacramental efficacy upon a minister's holiness, rejection of clerical authority and tithes, extraordinary claims for Tanchelm's own reception of the Holy Spirit, and a personal cult. Its sexual and ceremonial allegations use conventional heresiological motifs and cannot be treated as independently verified propositions.

\responsefield The canons asked Frederick to detain and judge Tanchelm and associates; the surviving letter is an accusation, not a preserved synodal sentence. Norbert later preached at Antwerp in defense of Catholic ministry and Eucharistic communion. No extant papal or ecumenical act names Tanchelm, and reports of a Cologne proceeding do not preserve an authoritative proposition list.

\aftermathfield Tanchelm was killed by an individual cleric or attendant, not executed under an extant conciliar sentence. His following soon disappears from identifiable evidence. Later references to Tanchelmites repeat the same accusatory dossier; the surviving evidence identifies neither an organized church nor a genetic link to other medieval dissenters.
\end{heresydossier}

\begin{heresydossier}{Petrobrusians and Henricians}{dossier:petrobrusians-henricians}
\objectfield Peter of Bruys preached in southern France in the early twelfth century; Henry of Lausanne led a related but not identical reform movement. Neither left an extant doctrinal book. Peter the Venerable's \emph{Against the Petrobrusians}, Bernard of Clairvaux, and ecclesiastical reports are principal witnesses. Evidence grade C.

\propositionfield Peter the Venerable attributes rejection of infant baptism, church buildings and altars, veneration of the cross, the Eucharistic sacrifice and real presence, and prayers or offerings for the dead. Henry's preaching stressed clerical corruption and unauthorized ministry; the degree to which he shared Peter's fivefold program remains uncertain.

\responsefield Local bishops and preachers opposed both movements. Lateran II canon 23 in 1139 anathematized those who condemned the Eucharist, infant baptism, priesthood and sacred orders, or lawful marriage; it provides an authoritative doctrinal boundary but does not name Peter or Henry or prove that each follower held every item. Bernard's mission answered Henry pastorally and sacramentally.

\aftermathfield Peter was killed by a crowd, not executed by a council; Henry was arrested and disappears from secure evidence. Later medieval polemic sometimes made them ancestors of unrelated dissent. With no movement archive, direct descent into Waldensian or Cathar communities cannot be presumed.
\end{heresydossier}

\begin{heresydossier}{Peter Abelard's censured propositions}{dossier:abelard-propositions}
\objectfield Peter Abelard's theological and ethical works were examined at Soissons in 1121 and again at Sens, whose disputed date is conventionally given as 1140 or 1141. Abelard's writings survive, but the capitula prepared by opponents excerpt and combine them. Evidence grade A/B for the texts and B for the exact judicial construction.

\propositionfield The Sens dossier raised claims about the divine persons and attributes, an identification of the Holy Spirit with the world-soul, limits upon what God can do, inherited guilt, the purpose of the Incarnation, and whether culpability lies solely in conscious contempt. The list must be checked against Abelard's context; it was not a confession of a school, and the later papal act did not assign a separate theological grade to each item.

\responsefield Soissons required Abelard to burn his first Trinitarian book and imposed monastic custody. At Sens he appealed before the full hearing concluded. Two papal letters followed in July 1141: Innocent II's \emph{Testante apostolo} confirmed condemnation of the transmitted propositions and writings in the aggregate and imposed silence upon Abelard, while a separate mandate ordered Abelard and Arnold of Brescia confined and their erroneous books destroyed. The censure is \emph{in globo}, not a warrant to call every extracted sentence formally heretical.

\aftermathfield Peter the Venerable received Abelard at Cluny and mediated a reconciliation with Bernard of Clairvaux; Abelard died as a monk in 1142. His logical and theological work remained influential. No Abelardian sect survived, and his submission belongs to the record as essentially as the censure.
\end{heresydossier}

\begin{heresydossier}{Arnold of Brescia and the Arnoldists}{dossier:arnoldists}
\objectfield Arnold of Brescia combined a reforming attack on clerical wealth with political action in the Roman commune. His own writings are lost; Bernard of Clairvaux, John of Salisbury, Otto of Freising, and papal records are interested witnesses. ``Arnoldists'' later appears as a group name in \emph{Ad abolendam}. Evidence grade C.

\propositionfield Opponents attribute the claim that clergy owning property, bishops exercising temporal lordship, and perhaps monks possessing goods could not validly or worthily exercise sacred ministry. The precise step from a moral demand for apostolic poverty to denial of sacramental authority is difficult to recover from the surviving polemic.

\responsefield Arnold was condemned in connection with Lateran II in 1139 and later excommunicated; Eugenius III and Adrian IV opposed his Roman program. \emph{Ad abolendam} in 1184 named Arnoldists among groups subject to anathema. Catholic reform distinguished the need to correct clerical avarice from a theory that temporal possession by itself extinguishes orders or sacramental power.

\aftermathfield Frederick Barbarossa's officials captured Arnold, and Roman civil authority executed and burned him in 1155; the death was not a conciliar sentence. Later use of his name for anti-clerical groups does not establish a continuous organized Arnoldist church.
\end{heresydossier}

\begin{heresydossier}{Gilbert of Poitiers at Reims}{dossier:gilbert-poitiers}
\objectfield Gilbert, bishop of Poitiers and commentator on Boethius, answered doctrinal accusations before Eugene III at Paris in 1147 and Reims in 1148. Gilbert's commentaries survive, but the proceeding is reconstructed from Otto of Freising, John of Salisbury, Geoffrey of Auxerre, and other reports that differ about wording and outcome. Evidence grade B/C.

\propositionfield The reported issues concern whether a real distinction between God and divinity makes the divine essence something other than God; how the one essence is predicated of the three persons; whether divine attributes are identical with God; and in what sense the divine nature was incarnate. These technical questions cannot be reduced safely to the slogan that Gilbert taught four gods.

\responsefield Eugene III and the cardinals required formulations preserving the identity of God and the one simple divine essence, the unity of the three persons in that essence, and the true Incarnation of the Son. No signed conciliar canons or papal bull of condemnation survives. Gilbert accepted the required corrections, amended disputed language, and retained his see; reports do not support a final personal verdict of obstinate heresy.

\aftermathfield Gilbert's school continued to shape scholastic semantics and Trinitarian analysis within Catholic theology. The Reims settlement illustrates correction without creation of a lasting named heresy: later technical distinctions must be assessed on their own meaning, not called Gilbertine by association.
\end{heresydossier}

\begin{heresydossier}{Cathars or Albigensians: diverse western dualisms}{dossier:cathars}
\objectfield From the mid-twelfth century, western sources describe dualist teachers and ritual communities under names including Cathars, Albigensians, Patarenes, and ``good Christians.'' The \emph{Book of the Two Principles} and ritual texts preserve some insiders' voices, while sermons, polemics, and inquisitorial depositions dominate other regions. Evidence ranges from A to C; modern scholars dispute the cohesion of a trans-European Cathar church.

\propositionfield Documented groups distinguished the good spiritual God from the maker or ruler of the visible world, rejected Catholic baptism and Eucharist, and treated the \emph{consolamentum} as the decisive rite. Some taught two eternal principles; others made the evil principle derivative. Ascetical rejection of marriage and animal food belonged especially to the perfected, not necessarily every hearer.

\responsefield Lateran III canon 27 in 1179 and Lucius III's \emph{Ad abolendam} in 1184 condemned named dualist groups and their protectors. Lateran IV constitution 1 in 1215 answered with one Creator of visible and invisible things, the goodness of created nature, one incarnate Christ, and sacramental baptism and Eucharist; constitution 3 established procedures and penalties against heretics. Dominican and other preaching missions supplied a pastoral and theological response.

\aftermathfield The Albigensian Crusade from 1209, confiscations, inquisitorial imprisonment, and executions were military, civil, and penal responses distinct from the creed. Organized dualist communities declined by the fourteenth century. Neither the evidence nor genetic metaphor justifies identifying modern dissenters as a surviving Cathar church.
\end{heresydossier}

\begin{heresydossier}{Waldensians or the Poor of Lyons}{dossier:waldensians}
\objectfield Valdes of Lyons embraced apostolic poverty and vernacular preaching in the 1170s. His followers first sought ecclesiastical recognition; reports of Lateran III, \emph{Ad abolendam}, professions, and later Waldensian texts show development rather than one founding doctrinal manifesto. Evidence grade A/B.

\propositionfield The original conflict concerned lay preaching without episcopal authorization. Later branches rejected oaths, capital punishment, purgatory, indulgences, and aspects of sacramental and hierarchical authority, but those later positions cannot all be attributed to Valdes at the beginning.

\responsefield Alexander III reportedly praised the poverty ideal while withholding authorization to preach without local approval. \emph{Ad abolendam} in 1184 excommunicated the Poor of Lyons for unauthorized preaching among other named groups. Innocent III received Durand of Huesca and others through a detailed profession of faith in 1208 and approved Catholic Poor communities, showing that reconciliation and ordered preaching were real alternatives to suppression.

\aftermathfield Some communities reconciled; others developed separate ministries and survived persecution into the Reformation, when they entered Reformed Protestant alliances. Civil banishment and executions varied by territory. Present Waldensian churches should be described through their current confession, not solely by a twelfth-century disciplinary sentence.
\end{heresydossier}

\begin{heresydossier}{The Humiliati: condemnation followed by reconciliation}{dossier:humiliati}
\objectfield The Humiliati began as lay penitential and working communities in northern Italy. Their organization contained different clerical, communal, and household forms; ``Humiliati'' in a prohibition does not denote an unchanging doctrinal sect. Papal letters and approved rules permit unusually clear institutional comparison. Evidence grade A/B.

\propositionfield The initial dispute concerned unauthorized assemblies and preaching rather than a surviving anti-sacramental creed. Some members continued to preach after prohibition; others sought a lawful form of common life and retained Catholic faith and sacraments.

\responsefield \emph{Ad abolendam} in 1184 censured those Humiliati or Poor of Lyons who presumed to preach without authorization. Under Innocent III, reconciled Humiliati received approved forms of life in 1201, with preaching and common discipline regulated under ecclesial authority. The two acts target different institutional states and cannot honestly be summarized as ``the Church declared the Humiliati heretical forever.''

\aftermathfield The approved order became an important Catholic religious and economic institution before its suppression for disciplinary collapse in 1571. This history demonstrates why a survey must track submission, reform, and juridical identity rather than inherit a static catalogue label.
\end{heresydossier}

\begin{heresydossier}{Passagians or Passagines}{dossier:passagians}
\objectfield Passagians are named in late-twelfth-century legislation and described by the anti-heretical \emph{Manifestatio} associated with Bonacursus and a \emph{Summa against Heretics} once attributed to Praepositinus. No target-side text, leader, or institutional record survives. Evidence grade C/D.

\propositionfield The polemical witnesses attribute obligatory observance of Mosaic ceremonial law by Christians, including circumcision, Sabbath, and food rules, together with a subordinationist account of Christ and denial of the received Trinitarian confession. The several items may describe one group, several groups, or a literary construct; \emph{Ad abolendam} itself states none of them.

\responsefield Lucius III's decretal \emph{Ad abolendam}, issued with the Verona assembly in 1184, placed the Passagians by name under anathema with other groups and established procedures against heresy. Later papal mandates repeated the canonical name. Those acts establish the classification and jurisdictional response, while the proposition content remains dependent upon the separate polemical tracts.

\aftermathfield No secure Passagian history can be followed beyond scattered legal repetition. Their attributed Mosaic observance must not be projected onto Jews, Jewish converts, or later Christian Sabbath observance, and repetition of a form letter is not evidence that a stable sect persisted in every named territory.
\end{heresydossier}

\begin{heresydossier}{Josephines: a name without recoverable propositions}{dossier:josephines}
\objectfield The Josephines or Josepini are securely known only because \emph{Ad abolendam} lists them between Passagians and Arnoldists. Later attempts to identify them with a Paulician leader named Joseph, a Cathar subgroup, or another medieval movement are conjectural. Evidence grade D.

\propositionfield No securely attributable Josephine proposition, book, leader, rite, or locality survives. The name therefore cannot support claims about dualism, Christology, sacraments, morals, or relation to Judaism. This dossier records the evidentiary absence rather than filling it from neighboring names in the decree.

\responsefield \emph{Ad abolendam} in 1184 subjected the named Josephines to anathema and the decree's procedural consequences. The act is an authoritative censure of a named object but contains no doctrinal specification by which a modern historian could identify an individual adherent.

\aftermathfield Canonical and heresiological lists sometimes repeated the name, but no independent aftermath is recoverable. Josephines remain in the census so that the legislative occurrence is not silently omitted; they cannot responsibly be narrated as a known heretical church.
\end{heresydossier}

\begin{heresydossier}{Amalricians and propositions known through a prosecution}{dossier:amalricians}
\objectfield Followers associated with the Paris master Amalric of B\`ene were investigated at Paris in 1210 after Amalric's death. Reports by Caesarius of Heisterbach, William the Breton, university and conciliar material preserve the case; no full Amalrician doctrinal book survives. Evidence grade B/C.

\propositionfield Sources attribute an undifferentiated presence of God in all things, an eschatological age of the Spirit in which ordinary sacraments and moral distinctions cease, and denial of bodily resurrection or final punishment. Which propositions belonged to Amalric, which to particular disciples, and which were judicial synthesis remains disputed.

\responsefield A provincial process at Paris condemned the accused; some recanted and received imprisonment, while unrepentant clerics were degraded and handed to secular power. Lateran IV constitution 2 in 1215 reprobated Amalric's teaching and described it as ``not so much heretical as insane,'' a deliberately distinctive censure rather than an exact proposition list.

\aftermathfield Secular authorities burned several convicted followers and exhumed Amalric's body; those acts were civil aftermath. ``Amaurian'' later became a loose label for pantheist or antinomian claims, so apparent recurrence requires new evidence.
\end{heresydossier}

\begin{heresydossier}{David of Dinant and the \emph{Quaternuli}}{dossier:david-dinant}
\objectfield David of Dinant wrote philosophical notebooks known as the \emph{Quaternuli}. The work is mostly lost, though fragments and summaries in Albert the Great, Thomas Aquinas, university records, and other scholastic witnesses permit a limited reconstruction. His relation to the Amalricians is unproved. Evidence grade B/C.

\propositionfield Opponents attribute an identity of God, prime matter, and intellect, so that all reality reduces to one underlying substance. Because much of the formulation comes from refuters applying their own metaphysical vocabulary, it is safer to identify the alleged identity proposition than to call David's whole philosophy a demonstrated pantheist system.

\responsefield A provincial synod at Paris in 1210 ordered David's \emph{Quaternuli} burned and prohibited their use. Robert of Cour\c{c}on's university statutes of 1215 renewed the reading prohibition while also regulating Aristotle's natural works. Lateran IV condemned Amalric's specified teaching but did not name David; the local academic acts must not be converted into an ecumenical sentence.

\aftermathfield David's personal fate is uncertain and no organized Davidian community is documented. The prohibitions stimulated refutation by later scholastics, but the eventual Catholic reception of Aristotle shows that a temporary university ban on books was not a condemnation of philosophy or natural inquiry as such.
\end{heresydossier}

\begin{heresydossier}{Joachim of Fiore's Trinitarian proposition}{dossier:joachim-fiore}
\objectfield Joachim of Fiore was an abbot and scriptural theologian who died in 1202 after submitting his works to Roman judgment. Lateran IV considered his lost treatise against Peter Lombard; later ``Joachite'' apocalyptic movements drew selectively on a wider corpus. Evidence grade B for the condemned text, A for the conciliar response.

\propositionfield Joachim accused Lombard of positing a fourth reality in God and described the unity of Father, Son, and Spirit through a collective analogy that the council judged insufficient to confess one simple divine essence. This specific Trinitarian error must be distinguished from Joachim's periodization of history and from later claims about an imminent spiritual age.

\responsefield Lateran IV constitution 2 defended Lombard's confession: each divine person is the one undivided divine reality, and no similarity between Creator and creature is so great that an even greater dissimilarity is absent. It condemned anyone defending the specified error but expressly declined to disparage the Florensian monastery, because Joachim had submitted his writings to the Apostolic See.

\aftermathfield Joachim was not personally sentenced as an obstinate heretic by the council. Later Spiritual Franciscans and apocalyptic writers adapted his schemes in ways requiring separate judgment; the condemnation of one Trinitarian argument is not a censure of every work or every use of a threefold historical pattern.
\end{heresydossier}

\begin{heresydossier}{The Eternal Gospel and radical Joachimite reception}{dossier:radical-joachimism}
\objectfield Gerard of Borgo San Donnino's lost \emph{Introductorius in Evangelium aeternum}, circulated at Paris in 1254 with works of Joachim of Fiore, transformed Joachim's historical exegesis into a more radical program. Its content survives principally in extracts assembled by opponents and the Anagni commission. Evidence grade B/C.

\propositionfield The extracted claims identify Joachim's books with an eternal gospel that would supersede the New Testament, place the decisive transition to the age of the Spirit about 1260, and imply that the existing Church's sacraments, hierarchy, or scriptural economy belonged to a passing age. These are Gerard's attributed conclusions and later Joachimite developments, not propositions condemned as Joachim's own in Lateran IV.

\responsefield A cardinal commission at Anagni examined and censured the \emph{Introductorius} in 1255 under Alexander IV, and its circulation was suppressed. The surviving protocol, rather than a universal papal constitution with a numbered censure, controls the account. Franciscan superiors removed Gerard from ministry and confined him; that internal penalty was distinct from the proposition-level examination.

\aftermathfield Apocalyptic expectations continued among some Spiritual Franciscans and through pseudonymous Joachimite works, while other readers used Joachim without the supersession claim. Gerard died in religious confinement. Neither Franciscan poverty nor a theology of historical development is by itself radical Joachimism.
\end{heresydossier}

\begin{heresydossier}{Paris 1270 and 1277: local censures of propositions}{dossier:paris-condemnations}
\objectfield Bishop Stephen Tempier of Paris censured thirteen propositions in 1270 and 219 propositions on 7 March 1277 amid disputes over Aristotelian philosophy, Latin Averroism, and the limits of theology. The list and preamble survive; attribution to Siger of Brabant, Boethius of Dacia, Thomas Aquinas, or an organized ``Averroist party'' varies by proposition. Evidence grade A/B.

\propositionfield The lists address one intellect for all humans, denial of personal immortality, cosmic necessity, eternity of the world, constraints upon divine omnipotence, and claims that philosophical demonstration can establish conclusions contrary to faith. They also contain technical theses with no proven defender and propositions whose meaning depends on scholastic context.

\responsefield Tempier prohibited teaching the listed propositions in his diocese and attached excommunication, acting as local bishop rather than an ecumenical council. Archbishop Robert Kilwardby issued a separate Oxford condemnation in 1277. The positive boundary protected creation, providence, divine freedom and omnipotence, personal intellect and responsibility, and the unity of truth.

\aftermathfield The censures redirected university debate but did not make all 219 articles dogmatic definitions or every Aristotelian a heretic. A Parisian act in 1325 revoked the condemnation insofar as it touched Thomas Aquinas. Historians continue to dispute the list's influence upon later science; simple causal triumphalism is unwarranted.
\end{heresydossier}

\begin{heresydossier}{Apostolic Brethren and Dolcinians}{dossier:apostolics-dolcinians}
\objectfield Gerard Segarelli founded the Apostolic Brethren near Parma about 1260 in imitation of apostolic poverty. After his execution in 1300, Fra Dolcino led a more militant apocalyptic remnant until 1307. Papal prohibitions, chronicles, and inquisitorial descriptions dominate; Dolcino's reported letters survive indirectly. Evidence grade B/C.

\propositionfield Early conflict concerned an unauthorized religious order and preaching after suppression. Later Dolcinian teaching joined radical poverty to an apocalyptic succession of church ages, denunciation of a corrupt hierarchy, and the claim that the remnant represented the true apostolic church. Judicial sources may systematize beliefs that changed under persecution.

\responsefield Lyons II constitution 23, \emph{Religionum diversitatem nimiam}, prohibited new religious orders without naming Segarelli's group. Honorius IV's \emph{Olim felicis recordationis} of 1286 applied that rule and ordered the Apostolics to disband; Nicholas IV renewed the command in 1290. Continued defiance brought inquisitorial prosecution and classification as heresy.

\aftermathfield Segarelli and later Dolcino and Margaret were executed by secular authorities after ecclesiastical proceedings; military forces also destroyed the Dolcinian refuge. Their deaths were not penalties written into Lyons II's religious-order canon. The later militant phase should not be projected wholesale onto every early penitent attracted to Segarelli.
\end{heresydossier}
